\scanpage{63}
% archive sheet number: 62
On notera en $=$ temps que \uline{par construction}, \deleted{\unsure{la construction|le constructeur}}
la formation de $\underline A^0_{X/S},A^1_{X/S},\psi_{X/S}$ commute au
changement de base quelconque.

[N.B. Il y a peut-être lieu de définir
\unsure{axiomatiquemt|axiomatiquement} un \uline{schém. d'Albanese} de $X/S$ par
la condition de représenter le foncteur $T_X(P)$, $P$
pr. hom. \unsure{ss|en|les} schém. abélien, muni d'une $X$-trivialisation,
et par le fait que la $=$ chose soit valable pour
le changt de base. Pour l'\uline{existence}, il \uline{suffit}
mais il n'est pas nécessaire sans doute que
l'on puisse trouver un sous-schém. abélien
de $\underline{\mathrm{Pic}}_{X/S}$ ayant $\underline{\mathrm{Pic}}^0_{X/S}$ comme ens sous-jacent.
Ainsi, il est plausible \unsure{que 'un|que un|qu'un} schém. d'Albanese
pour $X/S$ existe si $S$ est le spectre d'un corps
$k$, \deleted{\unsure{et}} si $\underline{\mathrm{Pic}}^0_{X/k}$ n'est pas propre sur $k$
[plus grand sous-schém. abélien de $\underline{\mathrm{Pic}}^0_{X/k}$ ; il faut
prouver évidemment une propriété universelle de
ce dernier, pour tt changement de base \ldots\ldots{} ]].

\diagram{$\left\{\vcenter{\hbox{\begin{tikzcd}[ampersand replacement=\&,row sep=small,column sep=small]
X\arrow[r] \& \underline A^1_{X/S}\\
\underline{\mathrm{Pic}}^{\unsure{0|\bullet}}_{X/S}\&\underline{\mathrm{Pic}}^0_{\underline A^1/S}\simeq\underline{\mathrm{Pic}}^0_{\underline A^0/S}\arrow[l]\arrow[dl,"\sim"']\\
\underline B\arrow[u]\&
\end{tikzcd}}}\right.$}
